% Task 10: the text of the task. Tasks.tex puts all tasks into one PDF; the code shown here is in Task10/.
\settitle{Task 10}{The decoder with two semaphores}

\begin{story}
Ticket VD-1503, from the e-mail of Tomasz in Task 9: ``Try to rewrite the buffer from VD-1482 with two semaphores:
one counts the free places, the other one the ready frames. I'd like to compare both versions at the next code
review.''
\end{story}

The template is the solution of Task 8, with condition variables. Rewrite it with two semaphores:

\begin{center}
\begin{tikzpicture}
  \node[actor=sIndigo, minimum width=2.4cm, minimum height=1.4cm, align=center, font=\small\bfseries] (dec) at (0,0)
    {\Large\faFilm\\[2pt] decoder};
  \foreach \k in {0,1,2} \node[cellused, minimum size=0.75cm, font=\scriptsize] at (2.5+\k*0.78,0) {\#\k};
  \foreach \k in {3,4} \node[cellfree, minimum size=0.75cm] at (2.5+\k*0.78,0) {};
  \draw[data] (dec.east) -- (2.08,0);
  \node[actor=sPink, minimum width=2.4cm, minimum height=1.4cm, align=center, font=\small\bfseries] (scr) at (7.8,0)
    {\Large\faTv\\[2pt] display};
  \draw[data] (6.04,0) -- (scr.west);
  \node[font=\scriptsize, text=nInk] at (4.06,0.65) {buffer, \texttt{CAPACITY} = 5 in this picture};
  % the two semaphores: counters with tokens
  \node[hw, minimum width=3.3cm, minimum height=1.05cm, align=left, font=\scriptsize] (fs) at (1.9,-1.65) {};
  \node[font=\scriptsize\ttfamily, text=nInk, anchor=north west] at (fs.north west) {free\_slots = 2};
  \foreach \k in {0,1} \fill[nInk] ($(fs.west)+(0.35+\k*0.4,-0.18)$) circle (0.13);
  \node[font=\scriptsize, text=nInk, anchor=north] at (fs.south) {how many free places};
  \node[hw, minimum width=3.3cm, minimum height=1.05cm, align=left, font=\scriptsize] (fl) at (6.0,-1.65) {};
  \node[font=\scriptsize\ttfamily, text=nInk, anchor=north west] at (fl.north west) {full\_slots = 3};
  \foreach \k in {0,1,2} \fill[nInk] ($(fl.west)+(0.35+\k*0.4,-0.18)$) circle (0.13);
  \node[font=\scriptsize, text=nInk, anchor=north] at (fl.south) {how many ready frames};
\end{tikzpicture}
\end{center}

\begin{itemize}
  \item \texttt{free\_slots}: how many free places are in the buffer;
  \item \texttt{full\_slots}: how many decoded frames wait in the buffer.
\end{itemize}
Keep \texttt{std::queue}. A semaphore of this size is declared as
\texttt{std::counting\_semaphore<CAPACITY> name\{start value\};} (header \texttt{<semaphore>}).

\codesection

The template is in \ghfolder{Task10}. This is \texttt{video.cpp}; \texttt{video.h} is the same as in Task 8.

\codefile[firstline=8]{video.cpp}

\runline

\section{Your tasks}

\begin{questions}
  \item What values do \texttt{free\_slots} and \texttt{full\_slots} have at the start? Why?
  \item Who calls \texttt{acquire()} and who calls \texttt{release()} on each of them? Write the new version.
  \item With semaphores, do you still need the mutex? What for?
  \item A colleague first locks the mutex and only then calls \texttt{acquire()} on a semaphore.
        Write this version and run it. What happens? Why? (Draw who waits for whom.)
  \item How do you now tell the display that the video is over? Why does the flag \texttt{finished}
        from Task 8 not work so easily here?
  \item Compare your version with the one with condition variables. Which one is shorter? Which one is easier to
        read? When would you choose a condition variable, and when a semaphore?
\end{questions}
